\subsection{Criterios de M-regularidad.}
\label{subsection:IV.19.5}

Retomamos aquí, completándolos, los criterios para que una sucesión de elementos de un anillo $A$ sea M-regular o M-cuasirregular (siendo $M$ un $A$-módulo), ya estudiados en ($0$, 15.1).

\begin{theorem}[19.5.1]
\label{IV.19.5.1}
Sean $A$ un anillo, $\mathbf f=(f_i)_{1\leq i\leq n}$ una sucesión de elementos de $A$ y $M$ un $A$-módulo. Consideremos las propiedades siguientes:
\begin{enumerate}
  \item[a)] La sucesión $\mathbf f$ es M-regular.
  \item[b)] Se tiene $\operatorname{H}_i(\mathbf f,M)=0$ para todo $i>0$ (III, 1.1.3).
  \item[b$'$)] Se tiene $\operatorname{H}_1(\mathbf f,M)=0$.
  \item[c)] La sucesión $\mathbf f$ es M-cuasirregular.
\end{enumerate}
Entonces se tienen las implicaciones
\[
  a)\Rightarrow b)\Rightarrow b')\qquad\text{et}\qquad a)\Rightarrow c).
\]
Pongamos $\mathfrak J=\sum_{i=1}^n f_iA$. Si los módulos $M_i=M/(\sum_{j=1}^i f_jM)$ son separados para la topología $\mathfrak J$-preádica (resp. si todo módulo cociente de un submódulo de $M$ es separado para la topología $\mathfrak J$-preádica), se tiene también la implicación $c)\Rightarrow a)$ (resp. $b')\Rightarrow a)$.
\end{theorem}

Las implicaciones $a)\Rightarrow c)$ y $c)\Rightarrow a)$ cuando los $M_i$ son separados para la topología $\mathfrak J$-preádica ya se han demostrado ($0$, 15.1.9), y se recuerdan únicamente por completitud. También se ha demostrado (III, 1.1.4 y 1.1.3.3) que $a)$ implica $b)$, y $b)$ implica trivialmente $b')$. Queda, pues, por ver que, cuando todo módulo cociente de un submódulo de $M$ es separado para la topología $\mathfrak J$-preádica, $b')$ implica $a)$. Razonemos por inducción sobre $n$. Para $n=1$, la afirmación se deduce trivialmente de las definiciones, pues entonces $\operatorname{H}_1(\mathbf f,M)$ no es sino el núcleo de la homotecia $z\mapsto f_1z$ en $M$ (III, 1.1.1 y 1.1.2). Supongamos, pues, que $n\geq2$, y pongamos $\mathbf f'=(f_i)_{1\leq i\leq n-1}$; con las notaciones de (III, 1.1.2), se tiene $\operatorname{K}_\bullet(\mathbf f,M)=\operatorname{K}_\bullet(f_n)\otimes_A\operatorname{K}_\bullet(\mathbf f',M)$. Se tiene entonces (III, 1.1.4.1) la sucesión exacta
\begin{equation}
  0\longrightarrow\operatorname{H}_0(f_n,\operatorname{H}_1(\mathbf f',M))\longrightarrow\operatorname{H}_1(\mathbf f,M)\longrightarrow\operatorname{H}_1(f_n,\operatorname{H}_0(\mathbf f',M))\longrightarrow0.
  \tag{19.5.1.1}\label{IV.19.5.1.1}
\end{equation}

\oldpage[IV-4]{205}

La relación $\operatorname{H}_1(\mathbf f,M)=0$ implica, pues, $\operatorname{H}_0(f_n,\operatorname{H}_1(\mathbf f',M))=0$ y $\operatorname{H}_1(f_n,\operatorname{H}_0(\mathbf f',M))=0$. La primera de estas dos relaciones significa que $f_n\operatorname{H}_1(\mathbf f',M)=\operatorname{H}_1(\mathbf f',M)$ (III, 1.1.3.5). Ahora bien, por definición (III, 1.1.1), $\operatorname{H}_1(\mathbf f',M)$ es isomorfo a un cociente de un submódulo $N$ de $M^{n-1}$; tomando sobre $M^{n-1}$ la filtración por los $M^j$ para $j\leq n-1$, sobre $N$ la filtración inducida y sobre $\operatorname{H}_1(\mathbf f',M)$ la filtración cociente de la de $N$, se obtiene sobre $\operatorname{H}_1(\mathbf f',M)$ una filtración finita cuyos cocientes son isomorfos a cocientes de submódulos de $M$, y por tanto son separados para la topología $\mathfrak J$-preádica por hipótesis. Como la relación $f_n\operatorname{H}_1(\mathbf f',M)=\operatorname{H}_1(\mathbf f',M)$ implica \emph{a fortiori} $\mathfrak J\operatorname{H}_1(\mathbf f',M)=\operatorname{H}_1(\mathbf f',M)$, se deduce de la hipótesis y de las observaciones anteriores que $\operatorname{H}_1(\mathbf f',M)=0$. La hipótesis de inducción demuestra entonces que la sucesión $\mathbf f'$ es M-regular. Por otra parte, la relación $\operatorname{H}_1(f_n,\operatorname{H}_0(\mathbf f',M))=0$ significa que $f_n$ es $(M/(\sum_{i=1}^{n-1}f_iM))$-regular; por tanto, la sucesión $\mathbf f$ es M-regular.

\begin{corollary}[19.5.2]
\label{IV.19.5.2}
Supongamos que $A$ es un anillo noetheriano, que los $f_i$ ($1\leq i\leq n$) pertenecen al radical de $A$ y que $M$ es un $A$-módulo de tipo finito. Entonces las cuatro condiciones $a)$, $b)$, $b')$, $c)$ de (19.5.1) son equivalentes.
\end{corollary}

En efecto, se sabe entonces ($0_{\mathrm I}$, 7.3.5) que todo $A$-módulo de tipo finito es separado para la topología $\mathfrak J$-preádica.

\begin{proposition}[19.5.3]
\label{IV.19.5.3}
Sean
\[
  0\longrightarrow M'\longrightarrow M\longrightarrow M''\longrightarrow0
\]
una sucesión exacta de $A$-módulos, $\mathbf f=(f_i)_{1\leq i\leq n}$ una sucesión M-regular y $\mathfrak J=\sum_{i=1}^n f_iA$. Consideremos las propiedades siguientes:
\begin{enumerate}
  \item[a)] La sucesión $\mathbf f$ es $M''$-regular.
  \item[b)] La filtración $\mathfrak J$-preádica de $M'$ está inducida sobre $M'$ por la filtración $\mathfrak J$-preádica de $M$ (en otras palabras, el homomorfismo canónico $\operatorname{gr}_{\mathfrak J}^*(M')\to\operatorname{gr}_{\mathfrak J}^*(M)$ es inyectivo).
  \item[c)] El homomorfismo canónico $M'/(\sum_{i=1}^n f_iM')\to M/(\sum_{i=1}^n f_iM)$ es inyectivo.
\end{enumerate}
Entonces se tienen las implicaciones $a)\Rightarrow b)\Rightarrow c)$, y $a)$ implica además que la sucesión $\mathbf f$ es $M'$-regular. Si todo cociente de un submódulo de $M''$ es separado para la topología $\mathfrak J$-preádica, las condiciones $a)$, $b)$ y $c)$ son equivalentes.
\end{proposition}

Es claro que $c)$ es consecuencia de $b)$. Demostremos que, bajo las hipótesis de la última afirmación, $c)$ implica $a)$: puesto que $\mathbf f$ es M-regular, se tiene, por (19.5.1), $\operatorname{H}_1(\mathbf f,M)=0$, de donde se obtiene una sucesión exacta, porción de la sucesión exacta de homología
\[
  0\longrightarrow\operatorname{H}_1(\mathbf f,M'')\longrightarrow\operatorname{H}_0(\mathbf f,M')\longrightarrow\operatorname{H}_0(\mathbf f,M).
\]
Ahora bien, la condición $c)$ expresa que el homomorfismo $\operatorname{H}_0(\mathbf f,M')\to\operatorname{H}_0(\mathbf f,M)$ es inyectivo (III, 1.1.3.5); equivale, pues, a $\operatorname{H}_1(\mathbf f,M'')=0$, de donde, en virtud de la hipótesis de separación y de (19.5.1), se sigue que $\mathbf f$ es $M''$-regular.

Demostremos a continuación que $a)$ implica $c)$ y que $\mathbf f$ es $M'$-regular, procediendo por inducción sobre $n$. Para $n=1$, como $f_1$ es M-regular, también es $M'$-regular, y la propiedad $c)$ no es sino el lema (3.4.1.4). Para $n>1$, la hipótesis de inducción muestra que

\oldpage[IV-4]{206}

si se pone $\mathbf f'=(f_1,\ldots,f_{n-1})$, la sucesión $\mathbf f'$ es $M'$-regular y se tiene, en virtud de $c)$ aplicada a $\mathbf f'$, una sucesión exacta
\[
  0\longrightarrow M'/(\textstyle\sum_{i=1}^{n-1}f_iM')\longrightarrow M/(\textstyle\sum_{i=1}^{n-1}f_iM)\longrightarrow M''/(\textstyle\sum_{i=1}^{n-1}f_iM'')\longrightarrow0.
\]
Por hipótesis, $f_n$ es $(M/(\sum_{i=1}^{n-1}f_iM))$-regular y $(M''/(\sum_{i=1}^{n-1}f_iM''))$-regular; el mismo razonamiento muestra, por una parte, que $f_n$ es $(M'/(\sum_{i=1}^{n-1}f_iM'))$-regular y, por otra, en virtud de (3.4.1.4), que la sucesión
\[
  0\longrightarrow M'/(\textstyle\sum_{i=1}^{n}f_iM')\longrightarrow M/(\textstyle\sum_{i=1}^{n}f_iM)\longrightarrow M''/(\textstyle\sum_{i=1}^{n}f_iM'')\longrightarrow0
\]
es exacta, de donde se sigue $c)$.

Demostremos, por último, que $a)$ implica $b)$. Como la sucesión $\mathbf f$ es entonces M-regular y $M'$-regular, y por tanto M-cuasirregular y $M'$-cuasirregular, se tienen isomorfismos canónicos
\[
  \operatorname{gr}_{\mathfrak J}^*(M')\xrightarrow{\sim}\operatorname{gr}_{\mathfrak J}^0(M')[T_1,\ldots,T_n],
  \qquad
  \operatorname{gr}_{\mathfrak J}^*(M)\xrightarrow{\sim}\operatorname{gr}_{\mathfrak J}^0(M)[T_1,\ldots,T_n]
\]
($0$, 15.1.7), y como se ha visto anteriormente que $\operatorname{gr}_{\mathfrak J}^0(M')\to\operatorname{gr}_{\mathfrak J}^0(M)$ es inyectivo, lo mismo ocurre con $\operatorname{gr}_{\mathfrak J}^*(M')\to\operatorname{gr}_{\mathfrak J}^*(M)$.

Obsérvese también que la equivalencia de las condiciones $a)$, $b)$, $c)$ de (19.5.3) es válida, en particular, cuando $A$ es noetheriano, $M''$ es un $A$-módulo de tipo finito y los $f_i$ pertenecen al radical de $A$.

\begin{env}[19.5.4]
\label{IV.19.5.4}
En lo que sigue de este número, conservamos las notaciones $A$, $\mathbf f$, $M$, $\mathfrak J$ de (19.5.1), y suponemos además que $M$ está provisto de una filtración decreciente $(M_k)$ formada por sub-$A$-módulos, tal que $M_0=M$. Recordemos ($0$, 15.1.5) que se define entonces sobre $M$ una segunda filtración decreciente formada por los submódulos
\[
  M_k'=M_k+\mathfrak JM_{k-1}+\cdots+\mathfrak J^{k-1}M_1+\mathfrak J^kM_0
\]
y que, si $\operatorname{gr}_\bullet(M)$ y $\operatorname{gr}_\bullet'(M)$ son los $A$-módulos graduados asociados a las filtraciones $(M_k)$ y $(M_k')$, respectivamente, se define un homomorfismo graduado sobreyectivo de grado $0$ ($0$, 15.1.5.2)
\[
  \psi_M:(\operatorname{gr}_\bullet(M)\otimes_A(A/\mathfrak J))[T_1,\ldots,T_n]\longrightarrow\operatorname{gr}_\bullet'(M).
\]
Recordemos también que se tiene el homomorfismo canónico sobreyectivo ($0$, 15.1.1.1)
\[
  \varphi_{\operatorname{gr}_\bullet(M)}:(\operatorname{gr}_\bullet(M)\otimes_A(A/\mathfrak J))[T_1,\ldots,T_n]\longrightarrow\operatorname{gr}_{\mathfrak J}^*(\operatorname{gr}_\bullet(M)).
\]
\end{env}

\begin{theorem}[19.5.5]
\label{IV.19.5.5}
Con las notaciones de (19.5.4), consideremos las propiedades siguientes:
\begin{enumerate}
  \item[a)] La sucesión $\mathbf f$ es $\operatorname{gr}_\bullet(M)$-regular.
  \item[b)] El homomorfismo canónico $\psi_M$ es biyectivo.
  \item[c)] El homomorfismo canónico $\varphi_{\operatorname{gr}_\bullet(M)}$ es biyectivo (en otras palabras ($0$, 15.1.7), la sucesión $\mathbf f$ es $\operatorname{gr}_\bullet(M)$-cuasirregular).
\end{enumerate}

\oldpage[IV-4]{207}

Se tienen entonces las implicaciones\footnote{La implicación $b)\Rightarrow c)$, sin hipótesis de separación sobre los $\operatorname{gr}_k(M)$, se debe a P. Deligne, cuya demostración reproducimos a continuación.}
\[
  a)\Rightarrow b)\Rightarrow c).
\]
Además, si, para todo entero $k\geq0$, los módulos cocientes de los submódulos de $\operatorname{gr}_k(M)$ son separados para la topología $\mathfrak J$-preádica (lo que ocurre cuando $A$ es noetheriano, los $f_i$ pertenecen al radical de $A$ y los $\operatorname{gr}_k(M)$ son $A$-módulos de tipo finito), las condiciones $a)$, $b)$ y $c)$ son equivalentes.
\end{theorem}

La implicación $a)\Rightarrow b)$ se ha demostrado en ($0$, 15.1.8); la implicación $c)\Rightarrow a)$ bajo la hipótesis de separación es un caso particular de ($0$, 15.1.9). Queda por demostrar que $b)$ implica $c)$, lo que se hará en varias etapas.

Denotaremos por $\mathfrak J$ el ideal $\sum_{i=1}^n f_iA$ engendrado por los $f_i$; para toda sucesión $\mathbf q=(q_1,\ldots,q_n)$ de $n$ enteros $\geq0$, pondremos $|\mathbf q|=\sum_{i=1}^nq_i$ y $\mathbf f^{\mathbf q}=f_1^{q_1}f_2^{q_2}\cdots f_n^{q_n}$.

\begin{env}[19.5.5.1]
\label{IV.19.5.5.1}
Para que la aplicación $\psi_M$ sea inyectiva (y por tanto biyectiva), es necesario y suficiente que, para $p\geq0$ y $q\geq0$, se verifique la condición siguiente:

$(\mathrm I_{q,p})$ Para toda familia $(x_{\mathbf q}^p)_{|\mathbf q|=q}$ de elementos de $M_p$, la relación
\[
  \sum_{|\mathbf q|=q}\mathbf f^{\mathbf q}x_{\mathbf q}^p\in\sum_{i=1}^q\mathfrak J^{q-i}M_{p+i}+M_{p+q+1}'
\]
implica $x_{\mathbf q}^p\in M_{p+1}+\mathfrak JM_p$ para todo $\mathbf q$ tal que $|\mathbf q|=q$.

Se procede como en ($0$, 15.1.6), considerando el sub-$A$-módulo $Q_k$ (resp. $Q_k'$) de los términos de grado $k$ del primer (resp. segundo) miembro de ($0$, 15.1.5.2). Se dota a $Q_k$ de la filtración
\[
  (Q_k)_i=\sum_{j\leq k-i}\left(\sum_{|\mathbf l|=j}(\operatorname{gr}_{k-j}(M)\otimes_A(A/\mathfrak J))T^{\mathbf l}\right)
\]
y a $Q_k'$ de la filtración imagen formada por los $\psi_M((Q_k)_i)=(Q_k')_i$; basta de nuevo demostrar que los homomorfismos $(\psi_M)_{ki}:\operatorname{gr}_i(Q_k)\to\operatorname{gr}_i(Q_k')$ son inyectivos. Ahora bien, se tiene
\[
  \operatorname{gr}_i(Q_k)=\sum_{|\mathbf l|=k-i}((M_i/M_{i+1})\otimes_A(A/\mathfrak J))T^{\mathbf l}
  =\sum_{|\mathbf l|=k-i}(M_i/(\mathfrak JM_i+M_{i+1}))T^{\mathbf l};
\]

\oldpage[IV-4]{208}

por otra parte, $(Q_k')_{i+1}$ es la imagen de $\sum_{h=0}^{k-i-1}\mathfrak J^{k-i-h-1}M_{i+h+1}$ en $M_k'/M_{k+1}'$. Decir que $(\psi_M)_{ki}$ es inyectivo equivale, pues, a formular la condición $(\mathrm I_{k-i,i})$; de donde se sigue nuestra afirmación.
\end{env}

\begin{env}[19.5.5.2]
\label{IV.19.5.5.2}
Para que la sucesión $\mathbf f$ sea $\operatorname{gr}_\bullet(M)$-cuasirregular, basta que, para $p\geq0$ y $q\geq0$, se verifique la condición siguiente:

$(\mathrm S_{q,p})$ Para toda familia $(x_{\mathbf q}^p)_{|\mathbf q|=q}$ de elementos de $M_p$, la relación
\[
  \sum_{|\mathbf q|=q}\mathbf f^{\mathbf q}x_{\mathbf q}^p\in M_{p+1}
\]
implica $x_{\mathbf q}^p\in M_{p+1}+\mathfrak JM_p$ para todo $\mathbf q$ tal que $|\mathbf q|=q$.

Por definición, decir que $\mathbf f$ es $\operatorname{gr}_\bullet(M)$-cuasirregular significa que, para todo $p\geq0$ y todo $q\geq0$, una relación
\[
  \sum_{|\mathbf q|=q}\mathbf f^{\mathbf q}x_{\mathbf q}^p\in M_{p+1}+\mathfrak J^{q+1}M_p
\]
para una familia $(x_{\mathbf q}^p)_{|\mathbf q|=q}$ de elementos de $M_p$ implica $x_{\mathbf q}^p\in M_{p+1}+\mathfrak JM_p$ para todo $\mathbf q$ tal que $|\mathbf q|=q$. Ahora bien, la hipótesis sobre la sucesión $(x_{\mathbf q}^p)$ significa que existe una familia $(y_{\mathbf q}^p)_{|\mathbf q|=q}$ de elementos de $\mathfrak JM_p$ tal que $\sum_{|\mathbf q|=q}\mathbf f^{\mathbf q}(x_{\mathbf q}^p-y_{\mathbf q}^p)\in M_{p+1}$. La condición $(\mathrm S_{q,p})$ implica entonces $x_{\mathbf q}^p-y_{\mathbf q}^p\in M_{p+1}+\mathfrak JM_p$ para todo $\mathbf q$ tal que $|\mathbf q|=q$, de donde $x_{\mathbf q}^p\in M_{p+1}+\mathfrak JM_p$; se ve, pues, que las condiciones $(\mathrm S_{q,p})$ implican que la sucesión $\mathbf f$ es $\operatorname{gr}_\bullet(M)$-cuasirregular.
\end{env}

Basta, pues, demostrar la proposición siguiente:

\begin{env}[19.5.5.3]
\label{IV.19.5.5.3}
Para todo $r\geq0$, si las condiciones $(\mathrm I_{q,p})$ se verifican para $p+q<r$, entonces las condiciones $(\mathrm S_{q,p})$ se verifican para $p+q<r$.

Introduzcamos las condiciones

$(\mathrm A_{q,p})$ Para toda familia $(x_{\mathbf q}^p)_{|\mathbf q|=q}$ de elementos de $M_p$, la relación
\[
  \sum_{|\mathbf q|=q}\mathbf f^{\mathbf q}x_{\mathbf q}^p\in M_{p+1}
\]
implica
\[
  \sum_{|\mathbf q|=q}\mathbf f^{\mathbf q}x_{\mathbf q}^p\in\sum_{i=1}^q\mathfrak J^{q-i}M_{p+i}.
\]

Es claro que la conjunción de $(\mathrm A_{q,p})$ y $(\mathrm I_{q,p})$ implica $(\mathrm S_{q,p})$. Por otra parte, las condiciones $(\mathrm A_{1,p})$ y $(\mathrm S_{0,p})$ son triviales. Por consiguiente, (19.5.5.3) se deducirá de la proposición siguiente:
\end{env}

\begin{env}[19.5.5.4]
\label{IV.19.5.5.4}
Para cualesquiera $p\geq0$ y $q\geq1$, la condición $(\mathrm A_{q+1,p})$ es consecuencia de la conjunción de las condiciones $(\mathrm A_{q,p})$, $(\mathrm I_{q-1,p+1})$, $(\mathrm I_{q-2,p+2})$, $\ldots$, $(\mathrm I_{0,p+q})$.

En efecto, una vez demostrada esta proposición, las condiciones $(\mathrm I_{q,p})$ que se suponen verificadas para $p+q<r$ implicarán, para cada $p\leq r$, la condición $(\mathrm A_{q,p})$ para $1\leq q<r-p$, por inducción sobre $q$.

Demostremos, pues, (19.5.5.4). Observemos que la condición $(\mathrm A_{q,p})$ equivale también a
\end{env}

\begin{env}[19.5.5.5]
\label{IV.19.5.5.5}
\[
  \mathfrak J^qM_p\cap M_{p+1}\subset\sum_{i=1}^q\mathfrak J^{q-i}M_{p+i}.
\]

Sea, pues, $m\in\mathfrak J^{q+1}M_p\cap M_{p+1}\subset\mathfrak J^qM_p\cap M_{p+1}$. Aplicando la condición $(\mathrm A_{q,p})$, se ve que existe, para cada $i$ tal que $1\leq i\leq q$, una familia $(y_{\mathbf q}^{p+i})_{|\mathbf q|=q-i}$ de elementos de $M_{p+i}$ tal que
\[
  m=\sum_{i=1}^q\left(\sum_{|\mathbf q|=q-i}\mathbf f^{\mathbf q}y_{\mathbf q}^{p+i}\right).
\]
Supongamos que, para $1\leq i<j$ ($j\geq1$), se ha demostrado que
\[
  y_{\mathbf q}^{p+i}\in\mathfrak JM_{p+i}+M_{p+i+1}\qquad(\text{pour }|\mathbf q|=q-i).
\]
De aquí se deduce, por definición (19.5.4),
\[
  m-\sum_{i<j}\left(\sum_{|\mathbf q|=q-i}\mathbf f^{\mathbf q}y_{\mathbf q}^{p+i}\right)\in M_{p+q+1}'.
\]
Puesto que, por hipótesis, se verifica $(\mathrm I_{q-j,p+j})$, se deduce
\[
  y_{\mathbf q}^{p+j}\in\mathfrak JM_{p+j}+M_{p+j+1}
\]

\oldpage[IV-4]{209}

y, por inducción sobre $j$, se ve que esta condición se verifica para todo $j$ tal que $1\leq j\leq q$. Se tiene, pues,
\[
  m\in\sum_{i=1}^q\mathfrak J^{q-i}(\mathfrak JM_{p+i}+M_{p+i+1})=\sum_{i=1}^{q+1}\mathfrak J^{q-i+1}M_{p+i}
\]
y así se ha demostrado $(\mathrm A_{q+1,p})$ y terminado la demostración de (19.5.5).
\end{env}

\begin{remarks}[19.5.6]
\label{IV.19.5.6}
\begin{enumerate}
  \item[(i)] Los resultados de este número conservan un sentido y son válidos cuando, en lugar de un $A$-módulo $M$ y de elementos $f_i\in A$, se toma un objeto $M$ de una categoría abeliana cualquiera y $n$ endomorfismos $f_i$ de $M$ que conmutan dos a dos; las demostraciones se adaptan sin dificultad a este caso más general. La hipótesis de separación para un $A$-módulo $N$ respecto de la topología $\mathfrak J$-preádica debe sustituirse aquí por la hipótesis de que todo subobjeto de $N$ contenido en todos los $\sum_i f_i^r(N)$ ($r$ entero arbitrario) es necesariamente nulo. La misma observación se aplica a los resultados de 19.7.
  \item[(ii)] Supongamos que $A$ es un anillo \emph{graduado} en grados $\geq0$. Si $M$ (resp. cada $\operatorname{gr}_p(M)$) es un $A$-módulo graduado en grados acotados inferiormente y si los $f_i$ no contienen ninguna componente homogénea de grado $0$, la hipótesis de separación de ($0$, 15.1.9) se verifica \emph{ipso facto} para $M$ (resp. cada $\operatorname{gr}_p(M)$), y se ve, por tanto, que en este caso se tiene la implicación $c)\Rightarrow a)$ en (19.5.1) (resp. (19.5.5)).
  \item[(iii)] Recordemos que, sin hipótesis de separación, la implicación $c)\Rightarrow b)$ deja de ser válida incluso para $n=1$ ($0$, 15.1.12, (iii)).
\end{enumerate}
\end{remarks}
